\chapter{\MakeUppercase{Planteamiento del problema dinámico del tumor avascular}}
%\chapter{Planteamiento del problema dinámico del tumor avascular}
\thispagestyle{fancy}
\section{Planteamiento del modelo matemático del crecimiento dinámico tumoral}
Consideramos que el tumor ocupa un dominio acotado $\Omega\subset\mathbb{R}^2$ con frontera $\Gamma$  que cambia con el tiempo $t$ medido en minutos. Las variables de la dinámica del crecimiento tumoral son la concentración de nutrientes y la presión intracelular.

La concentración de nutrientes al interior del tumor se representa mediante la función $\sigma$ $\equiv$ $\sigma(x,y,t)$ que obedece a la siguiente ecuación de difusión:
\begin{equation}
	D_N \Delta \sigma - \gamma (\sigma -\sigma_B )-\delta_N \sigma=0 \text{ en } \Omega
\end{equation}
donde $D_N$ es el coeficiente de difusión, $- \gamma (\sigma -\sigma_B )$ modela la velocidad de transferencia de nutrientes desde los vasos sanguíneos,  $\delta_N \sigma$ representa el consumo de nutrientes  por parte de la células vivas en el tumor. 

Adicionalmente el tumor es modelado como un fluido incompresible de velocidad $u\equiv u(x,y,t)$ que se rige por la ecuación de continuidad
\begin{equation}
	\nabla \cdot u = \gamma_T \sigma - \delta _T\text{ en } \Omega
\end{equation}
donde el lado derecho de la ecuación expresa la velocidad de reproducción  de las células del tumor menos la velocidad de muerte celular.

La velocidad está relacionada con el gradiente de presión por medio de la ley de Darcy:
\begin{equation}
	u=-w_t \nabla p \text{ en } \Omega
\end{equation}
donde $p\equiv p(x,y,t)$ es la presión del fluido al interior del tumor y $w_T$ representa la movilidad celular, lo que nos permite escribir la siguiente ecuación para la presión 
\begin{equation}
	\Delta p = -\frac{\gamma_T}{w_T} \sigma + \frac{\delta _T}{w_T}\text{ en } \Omega
\end{equation}

\section{Condiciones de borde en la frontera libre}
Asumimos que la concentración de nutrientes $\sigma_{out}$ en la frontera del tumor es constante e igual a la concentración de nutriente en el exterior del tumor.
\begin{equation}
	\sigma |_{\Gamma}  = \sigma_{out}
\end{equation}

La presión satisface la relación de Young-Laplace en la frontera 
\begin{equation}
	p |_{\Gamma}  = \gamma \kappa
\end{equation}
donde $\gamma$ representa la tensión superficial y $\kappa$ es la curvatura de la frontera $\partial\Omega$.

La velocidad normal $V_n$, con la que se desplaza la frontera del tumor es la componente normal de la velocidad $u$ 
\begin{equation}
	V_n = 	u |_{\Gamma} \cdot n  = 
	-w_T	\nabla p |_{\Gamma} \cdot n 
\end{equation}
donde $n$ es la normal unitaria exterior a $\Gamma$.

\section{Modelo adimensional}

Siguiendo \citep{cristini2008nonlinear} introducimos las siguientes constantes
\begin{align}
	L_D & = \dfrac{ D_N^{1/2}}{ (\gamma_B + \delta_N)^{1/2}}\\
	\lambda_R &= \frac{w_T \gamma}{ L^3_D}
\end{align}
y variables adimensionales
\begin{align}
	\overline{x} & = \dfrac{ x}{L_D}\\
	\overline{y} & = \dfrac{ y}{L_D}\\
	\overline{t} & = \lambda_R t\\	
	\overline{\sigma} & = \dfrac{ \sigma}{\sigma_{out}}\\
	\overline{p} & = \dfrac{ L_D}{\gamma }p					
\end{align}
Así el modelo de crecimiento tumoral queda expresado en las siguientes ecuaciones adimensionales:
\begin{align}
	&	\Delta \overline{\sigma}- \overline{\sigma}+  \frac{\gamma_B\sigma_B}{\sigma_{out}(\gamma_B + \delta_N)}=0 \text{ en } \overline{\Omega}\\
	&	\Delta \overline{p} = -\frac{\gamma_T \sigma_{out}}{\lambda_R} \overline{\sigma} + \frac{\delta _T}{\lambda_R}\text{ en } \overline{\Omega}\\
	&	\overline{\sigma} |_{\partial \overline{\Omega}}  = 1\\
	&	\overline{p} |_{\partial \overline{\Omega}}  = \kappa \\
	&	\overline{V}_n = 	
	-	\nabla \overline{p} |_{\partial \overline{\Omega}} \cdot n 
\end{align}
donde $\overline{\Omega}$ es el dominio luego del cambio de coordenadas con las variables adimensionales $(\overline{x},\overline{y})$.
\section{Modelo desacoplado}
El modelo acoplado anterior puede simplificarse si hacemos 
\begin{equation}
	\lambda_M= \gamma_T \sigma_{out}
\end{equation}
\begin{equation}
	B=\frac{\gamma_B\sigma_B}{\sigma_{out}(\gamma_B + \delta_N)}
\end{equation}
\begin{equation} 
	G=\frac{\lambda_M}{\lambda_R}(1-B)
	\label{eqn:G}
\end{equation}
\begin{equation}
	A=\frac{\dfrac{\delta_T}{\lambda_M}-B}{1-B}
	\label{eqn:A}
\end{equation}
\begin{equation}
	U=\frac{ \overline{\sigma}-B }{1-B}
\end{equation}
\begin{equation}
	P=\overline{p} + (1-U)G- AG\left( \frac{\overline{x}^2+\overline{y}^2}{4}\right)
\end{equation}
de modo que las nuevas variables $U$, $P$, $V_n$ satisfacen el siguiente sistema de ecuaciones:
\begin{align}
	-\Delta U + U  &= 0 \text{ en }  \overline{\Omega}\\
	U &= 1 	 \text{ en }  \partial \overline{\Omega}\\
	\Delta P  &= 0 \text{ en }  \overline{\Omega}\\
	P &= \kappa - AG\left( \frac{\overline{x}^2+\overline{y}^2}{4} \right)	 \text{ en }  \partial \overline{\Omega}\\
	V_n& =-\nabla P \cdot n + G \nabla U \cdot n- AG \frac{n\cdot (\overline{x},\overline{y})}{2} \text{ en }  \partial \overline{\Omega} \label{eqn:velocidadVn}
\end{align}	
Para simplificar la notación se omitirán  barras al escribir los dominios y cantidades relacionadas a las variables adimensionales.

\section{Evolución de la frontera libre}
Consideramos que $\Omega$ y $\Gamma$ pueden ser definidas mediante cierta función $\phi$ llamada también \emph{función conjunto de nivel}:
\[
\Omega = \{ \phi<0\}\quad y \quad  \Gamma
=\{ \phi =0\}.
\]
además suponemos que $\phi$ en $t=0$ se comporta como una función distancia dirigida 
\[
\phi(x)=
\left\{
\begin{aligned}
	-{\rm dist}(x,\Gamma)&,&x\in\Omega\\
	{\rm dist}(x,\Gamma)&,&x\not\in\Omega
\end{aligned}
\right.
\]
en tal caso $\phi$ es Lipschitz continua, y si es diferenciable en $x\not \in \Gamma$ entonces $|\nabla \phi|=1$ \citep{delfour1994shape}. Como $\Gamma$ es una  curva de nivel 0 para $\phi$ entonces la normal exterior a $\Gamma$ es
\[
n = \frac{\nabla \phi}{|\nabla \phi|}
\]

Por ejemplo, para $\Omega=\{(x,y):x^2+y^2<R^2\}$ podemos elegir $\phi(x,y)=\sqrt{x^2+y^2}-R$ y en consecuencia  $\nabla\phi=(x/\sqrt{x^2+y^2},y/\sqrt{x^2+y^2})$ excepto en el origen, y la normal unitaria es $n=(x/R,y/R)$ en $\Gamma$.

En términos de las curvas de nivel, la curvatura está dada por
\[
\kappa = \nabla \cdot n = 
\nabla \cdot \left(\frac{\nabla \phi }{ |\nabla \phi |} \right)
=\frac{ \phi_{xx} \phi_{y}^2 - 2 \phi_{x} \phi_{y} \phi_{xy}
	+\phi_{yy} \phi_{x}^2  }{  ( \phi_{x}^2+ \phi_{y}^2   )^{1/2}}
\] 
de igual manera si $\Gamma$ se puede representar en forma polar $\left\{ (r,\theta),r\geq 0,r=f(\theta) \right\}$ entonces la función conjunto de nivel será $\phi(r,\theta) = r-f(\theta)$ y la curvatura está dada por
\[
\kappa =\frac{r^2 +2(f'(\theta))^2-rf''(\theta)}{
	\left[
	(f'(\theta))^2+r^2
	\right]^{3/2}
}
\]
y $
\nabla \phi =(\cos\theta,\sen\theta)+\frac{f'(\theta)}{r}(\sen\theta,-\cos\theta)$ de modo que $|\nabla \phi|=\sqrt{1+\frac{f'(\theta)^2}{r^2}}$.

En cualquier instante $t$ la frontera del tumor $\Gamma(t)$ está dada por el conjunto de nivel 0,
\[
\Gamma(t) =\left\{ (x,y): \phi(x,y,t)=0\right\}
\] y si una partícula permanece en la  posición $(x(t),y(t))$ sobre la frontera, entonces  $\phi(x(t),y(t),t)=0$, por lo tanto
\[
\begin{aligned}	
	\frac{d \phi}{dt}=0&= \frac{\partial \phi}{ \partial t}
	+\frac{\partial \phi}{ \partial x} \frac{d x}{dt}
	+\frac{\partial \phi}{ \partial y} \frac{d y}{dt}\\
	&=\frac{\partial \phi}{ \partial t}+\nabla \phi \cdot V
\end{aligned}
\]
donde $V$ es la velocidad de la partícula. Como 
la normal unitaria en la frontera está dada por $
n = \frac{\nabla \phi}{|\nabla \phi|}$ y $V=V_n n$  entonces
\[
\begin{aligned}	
	\frac{\partial \phi}{ \partial t}+V_n |\nabla \phi |=0\text{ en } \Gamma(t)
\end{aligned}
\]
esta ecuación solo es válida en la frontera, sin embargo al introducir la función conjunto de nivel podemos resolver la ecuación en un dominio fijo $\mathcal{O}\subset \mathbb{R}^{2}$, para ello es necesario encontrar una extensión $S=(S_1, S_2)$ para la velocidad $V$ de tal modo que 
\[
S_1 = V_n n_1, S_2 = V_n n_2\text{ en }\Gamma
\]
y así plantear la siguiente ecuación de advección para $\phi$
\begin{equation}
	\frac{\partial \phi}{ \partial t} + S \cdot \nabla \phi=0\text{ en } \mathcal{O}
	\label{eqn:evolucionconjuntodenivel}
\end{equation}
\section{Extensión de la velocidad de crecimiento}
Decimos que $u_{ext}$ es la extensión de $u_0$ definida en $\Gamma$ a todo el dominio fijo $\mathcal{O}$ si es la solución del siguiente problema:
\begin{equation}
	\begin{cases}
		\nabla \phi \cdot \nabla u =0& \text{ en }\mathcal{O}\\
		u =u_0 &\text{ en }\Gamma
	\end{cases}	
	\label{eqn:extension}
\end{equation}
Sea $S=(S_1, S_2)$ una extensión de la velocidad $V=(v_1,v_2)$ en la frontera del tumor, de modo que $S_1$ extiende $v_1$ y $S_2$ extiende $v_2$, es decir  $\nabla S_1 \cdot \nabla\phi =0$ y $\nabla S_2 \cdot \nabla\phi =0$ por lo tanto
\[
(S_1)_x \phi_x +(S_1)_y \phi_y=0
\]
\[
(S_2)_x \phi_x +(S_2)_y \phi_y=0
\]
Derivamos la expresión dada en la ecuación \ref{eqn:evolucionconjuntodenivel} respecto de $x$
\[
\frac{ \partial\phi_x}{ \partial t} + S_1 \phi_{xx}  + S_2 \phi_{xy}  +  (S_1)_x \phi_{x}  + (S_2)_x \phi_{y} =0
\]
 y también respecto de $y$ 
\[
\frac{\partial \phi_y}{ \partial t} + S_1 \phi_{xy}  + S_2 \phi_{yy}  +  (S_1)_y \phi_{x}  + (S_2)_y \phi_{y} =0
\]
multiplicamos por $\phi_x$, $\phi_y$ respectivamente y sumamos
\[
\frac{\partial |\nabla \phi|^2}{ \partial t} + 
\phi_x(\nabla S_1 \cdot \nabla\phi ) + \phi_y(\nabla S_2 \cdot \nabla\phi )
+ \frac{1}{2} S\cdot \nabla (\|\nabla \phi\|^2)=0
\]
de modo que si $\|\nabla \phi\|^2=1$ entonces $\frac{\partial |\nabla \phi|^2}{ \partial t}=0$, es decir la propiedad de la distancia dirigida es conservada.

\begin{defi}
Sea $F$ una función de $V$ en $\mathbb{R}$. Si existe
\[
\lim_{t\to 0^+} \frac{F(u+tv)-F(u)}{t}
\]
entonces será llamado derivada direccional de $F$ en la dirección $v$ y denotado por $F'(u;v)$. Si $F'(u;v)$ existe para todo $v\in V$ entonces $F$ se dice G\^{a}teaux diferenciable en $u$ y la aplicación lineal $\partial F[u]:V\to\mathbb{R}$ tal que $\partial F[u](v)=F'(u;v)$ se llama derivada de G\^{a}teaux de $F$ en $u$.
\end{defi}

Para formular la ecuación \ref{eqn:extension} como un problema elíptico definimos 
\[
R(v)= \frac{1}{2}\int_\mathcal{O} (\nabla \phi \cdot \nabla v)^2 + \frac{\gamma}{2}\int_\Gamma (v-u_0)^2
\] de modo que si $u$ extiende $u_0$ entonces de acuerdo a \citep{utz2018high}  se resuelve el  siguiente problema de minimización: 
\[
\min_{v\in H^1(\mathcal{O})} R(v) 
\]
luego la derivada de G\^{a}teaux es
\[
\partial R[u](v)=\lim_{t\to 0}\frac{R(u+tv)-R(u)}{t} = \int_\mathcal{O} (\nabla \phi \cdot \nabla u)(\nabla \phi \cdot \nabla v)+
\gamma \int_\Gamma (u-u_0)v
\]
formalmente podemos multiplicar por $\nabla \phi \cdot \nabla v$ y luego 
\[
\int_\mathcal{O} (\nabla \phi \cdot \nabla u)(\nabla \phi \cdot \nabla v)=-\int_\mathcal{O}( \nabla \cdot (A \nabla u))v+\int_{\partial\mathcal{O}}    ((A \nabla u) \cdot n)v
\]
donde $A=\begin{bmatrix} \phi_x^2 &  \phi_{x}\phi_{y} \\ \phi_{x}\phi_{y} & \phi_y^2 \end{bmatrix}$, lo que nos lleva al siguiente problema elíptico
\begin{equation}
	\begin{aligned}
		\nabla \cdot (A \nabla u)=0 & \text{ en }  \mathcal{O}\setminus \Gamma \\
		(A \nabla u) \cdot n=0 &	 \text{ en }  \partial \mathcal{O}\\
		u=u_0 &	 \text{ en } \Gamma
	\end{aligned}
\end{equation}
que debe resolverse para encontrar $S_1$ y $S_2$.


\section{Solución circular}
Recordemos que la ecuación 
\[
-\Delta u + \beta u =0
\]
es equivalente en coordenadas polares a
\[
u_{rr}+\frac{1}{r^2}u_{\theta\theta}+\frac{1}{r}u_r-\beta u=0
\]
si consideramos el caso de las soluciones con  simetría radial, es decir $u_{\theta}=0$, y $\partial\Omega=\left\{(r;\theta): r=R\right\}$  entonces $\kappa=1/R$, por lo tanto la ecuación para la presión se reduce a:
\[
\begin{aligned}
	P''(r)+\frac{1}{r}P'(r) &= 0 \text{ en }  \left]0,R\right[\\
	P(R) &= \frac{1}{R} - AG \frac{R}{4}  \\
\end{aligned}
\]	
cuya solución no singular es constante: $P(r) = \frac{1}{R} - AG \frac{R}{4},r\in[0,R]$.

La ecuación para la concentración de nutrientes se reduce a la ecuación diferencial de segundo orden del tipo
\begin{equation}
	\begin{aligned}
		r^2 U''(r)+rU'(r)-r^2U(r)&= 0 \text{ en }  \left]0,R\right[\\
		U(R) &= 1\\
	\end{aligned}
	\label{eqn:edocircular}
\end{equation}
que tiene como solución no singular  $U(r)=\frac{I_0(r)}{I_0(R)},r\in[0,R]$, donde $I_0(x)$ es una función modificada de Bessel de primer tipo \citep{gradshteyn2014table} cuya forma explícita es: 
\[
I_p(x)=\left( \frac{x}{2} \right)^p \sum_{k=0}^\infty \frac{x^{2k}}{4^k k!\Gamma(p+k+1)},x\in\mathbb{R},p>-1
\]
en particular
\[
I_0(x)=\sum_{k=0}^\infty \frac{(\frac{1}{4}x^2)^k}{(k!)^2}
\]
luego la velocidad de crecimiento del radio del tumor $R$, será
\begin{equation}
	\begin{aligned}
		V=\frac{dR}{dt}&=G\nabla U\cdot n-AG\frac{n\cdot(x,y)}{2}\\
		&=GU'(R)-AG\frac{R}{2}\\
		&=G\frac{I_1(R)}{I_0(R)}-AG\frac{R}{2}	
	\end{aligned}
	\label{eqn:edoradiotumor}
\end{equation}
De \citep{simpson1984some} vemos que $w_0:\left]0,+\infty\right[\to\mathbb{R}$ definida por $w_0(x)=\dfrac{xI_0(x)}{I_1(x)}$ es estrictamente creciente y convexa, %además de \citep{baricz2010bounds} tenemos que
%\[
%\frac{I_1(x)}{xI_0(x)}=\sum_{n\geq 1}\frac{2}{x^2+j^2_{0,n}}
%\]eqn:phi
%donde $j_{0,n}$, 
además en \citep{garofalo2020two} se muestra que:
\[
\lim_{x\to 0^+} 2\frac{(p+1)}{x}\frac{I_{p+1}}{I_p}=1,\quad p>-1
\]
y por lo tanto para $p=0$ concluimos que \[
\lim_{x\to 0^+} w_0(x)=2\]

La siguiente desigualdad es obtenida en \citep{yang2017sharp}:
\[
\sqrt{x^2+(p-1)^2}+p-1<\frac{xI_{p-1}}{I_p}<\sqrt{x^2+p^2}+p,\quad p\geq 1,x>0
\]
luego 
\[
x<w_0<\sqrt{x^2+1}+1
\]
y por lo tanto
\[
\lim_{x\to +\infty} w_0(x)=+\infty	
\]
Podemos escribir la velocidad de crecimiento del tumor como
\[
V= GR\left(\frac{1}{w_0(R)} -\frac{A}{2}\right)
\]
y analizar los siguientes casos:
\begin{itemize}
	\item $G\geq 0$, $A>0$. 
	
	Si $0<A<1$ entonces existe un único $R_\infty>0 $ tal que $w_0(R_\infty)=\frac{2}{A}$, por lo que  $V(R_\infty)=0$ y por tanto $R=R_\infty$ es una solución estacionaria de la ecuación de crecimiento tumoral, además
	\[
	R_\infty<\frac{2}{A}<\sqrt{	R_\infty^2+1}+1\\
	\]
	o
	\begin{equation}
		\sqrt{(\frac{2}{A}-1)^2-1}<R_{\infty}<\frac{2}{A}
		\label{eqn:radioestacionario}
	\end{equation}
	Si $A \geq 1$ entonces $V<0$ es decir el tumor decrece y tiende a desaparecer.
	\item $A\leq 0$ y $G\geq 0$. Entonces 
	\[
	V=G\frac{I_1}{I_0}-AG\frac{R}{2}>0
	\]
	por lo que $R$ es creciente, es decir el tumor crece de manera acelerada.
	\item  $G<0$. El crecimiento es acelerado o alcanza un punto de equilibrio si $A>0$ de manera similar al primer caso. El tumor tiende a desaparecer si $A<0$ pues tendremos que $V<0$.
	
\end{itemize}
De acuerdo a \cite{hogea2005computational} los regímenes de crecimiento tumoral pueden ser caracterizados por los parámetros $G$ y $A$:
\begin{itemize}
	\item Régimen de baja vascularización: $G\geq 0$, $A>0$. Corresponde a poca presencia de vasos sanguíneos al interior del tumor y por lo tanto los nutrientes no son suficientes para acelerar el crecimiento tumoral. En el caso de un tumor con frontera circular si $0<A<1$ se alcanza una solución estacionaria (\ref{eqn:radioestacionario}), si $A\geq 1$ el tumor tiende a desaparecer. 
	\item Régimen de vascularización moderada: $G\geq 0$, $A\leq 0$. En esta etapa el aumento de vasos sanguíneos al interior del tumor brinda una mayor cantidad de nutrientes y por lo tanto el crecimiento del tumor se acelera.
	\item Régimen de alta vascularización: $G<0$. En esta etapa el tumor ha desarrollado una densa red de vasos sanguíneos esto puede llevar a un crecimiento acelerado del tumor o a su estabilización dependiendo  si $A>0$ o $A<0$.
\end{itemize}

%
%\section{Soluciones lineales}
%Consideramos que la frontera del tumor esta representada por
%\[
%r(t,\theta) = R(t) + \delta(t)\cos(l\theta)
%\]
%por lo tanto

%%%%%%%%%%%%%%%%%%%%%%%%%%%%%%%%%%%%%%%%%%


%%%%%%%%%%%%%%%%%%%%%%%%%%%%%%%%%%%%%%%%%%
%	
%	\begin{figure}[H]
	%		\centering
	%		\begin{tikzpicture}
		%			\begin{loglogaxis}[xlabel=$h$,ylabel=$\epsilon_h$,
			%				legend style=
			%				{
				%					at={(1,1)},
				%					anchor=north west,
				%					draw=none,
				%					fill=none
				%				},
			%				legend cell align=left,]
			%				\addplot[blue,line width=1pt,domain=0.04:3] {exp(1.3266*ln(x)-3.7013)};
			%				\addplot[red,line width=1pt,domain=0.04:3] {exp(1.0790*ln(x)-0.9396)};	
			%				\addplot[only marks,color=blue,mark=otimes*] coordinates {
				%					(2.82843 ,  0.0891674)
				%					(1.41421 ,  0.0436996)
				%					(0.707107 ,  0.0149696)
				%					(0.353553 , 0.00679741)
				%					(0.176777 ,  0.0022653)
				%					(0.0883883 , 0.00108145)
				%					(0.0441942 ,0.000369282)
				%				};
			%				\addplot[only marks,color=red,mark=square*] coordinates {
				%					(2.82843 ,   1.79489)
				%					(1.41421 ,  0.344859)
				%					(0.707107 ,   0.284266)
				%					(0.353553 ,0.141762)
				%					(0.176777 ,  0.0513602)
				%					(0.0883883 , 0.0245787)
				%					(0.0441942 ,0.0171774)
				%				};
			%				\legend
			%				{
				%					$r=4$,
				%					$r=2.1 + 0.5\cos(2\theta)$,
				%				}
			%			\end{loglogaxis}
		%		\end{tikzpicture}
	%		\caption{Error en la curvatura de una circunferencia de radio 4}
	%		\label{fig:errorcurvatura}
	%	\end{figure}
%	\begin{table}[h]
	%		\centering
	%		\begin{tabular}{ccc}
		%			$h$ & $H_k$ & $\epsilon_h$ \\
		%			\hline
		%			2.8284 &  0.269424 &  0.0892\\
		%			1.4142 & 0.259138  & 0.0437\\
		%			0.7071 &  0.252702 & 0.0150\\
		%			0.3536 &  0.250265 & 0.0068\\
		%			0.1768 &   0.250189 & 0.0023\\
		%			0.0884 &  0.25002  & 0.0011\\
		%			0.0442 &   0.250013 & 0.0004\\
		%		\end{tabular}
	%		\caption{}
	%		\label{tab:curvaturacirculo}
	%	\end{table}

%\[
%\begin{aligned}
%	-\Delta u + \alpha u  = 0 \text{ en }  \Omega\\
%	u = g 	 \text{ en }  \partial \Omega\\
%\end{aligned}
%\]	
%\section{Definición del dominio de estudio simétrico y no simétrico}
%\section{Condiciones de frontera libre del dominio físico del tumor avascular}
%\section{Espacios funcionales de estudio de las variables presión-velocidad del comportamiento avascular}
%\section{Formulación variacional del problema dinámico avascular}
